% Part:sets-functions-relations
% Chapter: sets
% Section: reduction

\documentclass[../../../include/open-logic-section]{subfiles}

\begin{document}

\olfileid{sfr}{siz}{red}

\olsection{راکمونه}

\begin{editorial}
  دا برخه په راکمونه د شمېر نۀ وړ والے ثابتوي، او پايلې ئې د
  \olref[nen]{sec} له پايلو سره سمون خوري۔ يوه بله، لږه لنډه
  بڼه چې د \olref[nen-alt]{sec} له پايلو سره سمون خوري، په
  \olref[red-alt]{sec} کښې ورکړل شوې ده۔
\end{editorial}

موږ په قطري استدلال وښودل چې $\Pow{\PosInt}$ !!{nonenumerable}
دے۔ له دې مخکښې مو ثابته کړې وه چې $\Bin^\omega$، يعنې د
$0$ګانو او $1$ګانو د ټولو نامتناهي تسلسلونو سټ،
!!{nonenumerable} دے۔ دا د $\Pow{\PosInt}$ د !!{nonenumerable}
کېدو د ثابتولو يوه بله لار ده: وښيئ چې \emph{که $\Pow{\PosInt}$
!!{enumerable} وي، نو $\Bin^\omega$ هم !!{enumerable} دے}۔ ځکه
پوهېږو چې $\Bin^\omega$ !!{enumerable} نۀ دے، $\Pow{\PosInt}$
هم داسې نۀ شي کېدے۔ دې ته د يوې مسئلې بلې ته \emph{راکمونه}
ويل کېږي؛ په دې حالت کښې د $\Bin^\omega$ د شمېرلو مسئله د
$\Pow{\PosInt}$ د شمېرلو مسئلې ته راکموو۔ د وروستۍ مسئلې يو حل،
يعنې د $\Pow{\PosInt}$ يوه شمېرنه، به د لومړۍ مسئلې حل، يعنې د
$\Bin^\omega$ يوه شمېرنه، راکړي۔

د يوه سټ~$B$ د شمېرلو مسئله څنګه د يوه سټ~$A$ د شمېرلو مسئلې
ته راکموو؟ د~$A$ د يوې شمېرنې د~$B$ په شمېرنه د اړولو يوه لار
ورکوو۔ د دې تر ټولو اسانه لار د يوې !!a{surjective} تابع
$f\colon A \to B$ تعريفول دي۔ که $x_1$، $x_2$، \dots{} د~$A$
شمېرنه وي، نو $f(x_1)$، $f(x_2)$، \dots{} به د~$B$ شمېرنه وي۔
زموږ په حالت کښې د يوې پر هدف سټ بشپړې تابع
$f\colon \Pow{\PosInt} \to \Bin^\omega$ په لټه کښې يو۔

\begin{prob}
وښيئ چې که يوه !!{injective} تابع $g\colon B \to A$ وي او
$B$~!!{nonenumerable} وي، نو $A$ هم همداسې دے۔ دا په ښودلو
وکړئ چې څنګه د~$g$ په کارولو د~$A$ يوه شمېرنه د~$B$ په شمېرنه
اړولے شئ۔
\end{prob}

\begin{proof}[په راکمونه د {\olref[nen]{thm:nonenum-pownat}} ثبوت]
فرض کړئ چې $\Pow{\PosInt}$ !!{enumerable} وے، او له دې امله د
هغه يوه شمېرنه $Z_{1}$، $Z_{2}$، $Z_{3}$، \dots وے۔

تابع $f \colon \Pow{\PosInt} \to \Bin^\omega$ داسې تعريف کړئ
چې $f(Z)$ هغه تسلسل $s$ وي چې $s(n) = 1$ هله او يوازې هله
چې $n \in Z$ وي، او په بل حالت کښې $s(n) = 0$ وي۔
\textbf{د ژباړې د نسخې سمون (OLSIZ-004):} په انګرېزي نثر کښې
د تسلسل بې تړاوه شاخص د خپل سري ورودي سټ له نښې سره بدل شو،
څو د تابع تعريف روښانه وي۔ دا په څرګند ډول تابع تعريفوي، ځکه
کله چې $Z \subseteq \PosInt$ وي، هر $n \in \PosInt$ يا د $Z$
يو !!a{element} وي يا نۀ وي۔ د بېلګې په ډول، د مثبتو جفتو
عددونو سټ $2\PosInt = \{2, 4, 6, \dots\}$ تسلسل
$010101\dots$ ته نقشېږي، تش سټ $0000\dots$ ته او پخپله سټ
$\PosInt$، $1111\dots$ ته نقشېږي۔

دا !!{surjective} هم ده: د $0$ګانو او $1$ګانو هر تسلسل د مثبتو
صحيح عددونو له يوه سټ سره سمون خوري، يعنې له هغه سټ سره چې غړي
ئې هغه صحيح عددونه دي کوم چې د تسلسل د~$1$ لرونکو مقامونو سره
سمون خوري۔ په لا دقيقه توګه، فرض کړئ $s \in \Bin^\omega$۔
$Z \subseteq \PosInt$ په لاندې ډول تعريف کړئ:
\[
Z = \Setabs{n \in \PosInt}{s(n) = 1}
\]
بيا $f(Z) = s$، لکه چې د~$f$ تعريف ته په کتلو تاييدولے شو۔

اوس دا لړ په پام کښې ونيسئ:
\[
f(Z_1), f(Z_2), f(Z_3), \dots
\]
ځکه چې $f$ !!{surjective} ده، د $\Bin^\omega$ هر غړے بايد د
کوم ورودي دپاره د~$f$ د قيمت په توګه راشي، او نو بايد په لړ کښې
راشي۔ له دې امله دا لړ بايد ټول~$\Bin^\omega$ وشمېري۔

نو که $\Pow{\PosInt}$ !!{enumerable} وے، $\Bin^\omega$ به
!!{enumerable} وے۔ خو $\Bin^\omega$ !!{nonenumerable} دے
(\olref[nen]{thm:nonenum-bin-omega})۔ له دې امله
$\Pow{\PosInt}$ !!{nonenumerable} دے۔
\end{proof}

\begin{explain}
د راکمونې په لوري کښې ګډوډېدل اسانه دي۔ د بېلګې په ډول، يوه
!!a{surjective} تابع $g \colon \Bin^\omega \to B$ دا \emph{نۀ}
ثابتوي چې $B$ !!{nonenumerable} دے۔ (تابع
$g \colon \Bin^\omega \to \Bin$ په پام کښې ونيسئ چې په
$g(s) = s(1)$ تعريف شوې؛ دا تابع د $0$ګانو او $1$ګانو يو تسلسل
د هغه لومړي !!{element} ته نقشوي۔ دا !!{surjective} ده، ځکه ځينې
تسلسلونه په $0$ او ځينې په $1$ پيلېږي۔ خو $\Bin$ متناهي دے۔)
دا هم په ياد ولرئ چې تابع~$f$ بايد !!{surjective} وي، ګنې
استدلال مخکښې نۀ ځي: بيا به دا ډاډ نۀ وي چې $f(x_1)$، $f(x_2)$،
\dots{} د~$B$ ټول !!{element}s لري۔ د بېلګې په ډول،
\[
h(n) = \underbrace{000\dots0}_{\text{$n$ $0$ګانې}}111\dots
\]
تابع $h\colon \PosInt \to \Bin^\omega$ تعريفوي، خو $\PosInt$
!!{enumerable} دے۔
\textbf{د ژباړې د نسخې سمون (OLSIZ-005):} اصلي ښودل شوے حاصل
متناهي توريز تسلسل و، نو د يادې تابع د هدف غړے نۀ و۔ دلته تر
لومړنيو صفرونو وروسته د يوونو نامتناهي لکۍ ورزياته شوې؛ د راکمونې
د لوري په اړه پايله هماغه پاتې کېږي۔
\end{explain}

\begin{prob}
په راکمونې استدلال وښيئ چې د مثبتو صحيح عددونو د جوړو د
\emph{سټونو} سټ !!{nonenumerable} دے۔
\end{prob}

\begin{prob}\label{sfr:siz:red:prob:nat-nat}
  په راکمونې استدلال وښيئ چې د ټولو تابعو $f\colon \Nat \to \Nat$
  سټ~$X$ !!{nonenumerable} دے۔ (لارښوونه: له $X$ څخه
  ~$\Bin^\omega$ ته يوه پر هدف سټ بشپړه تابع ورکړئ۔)
\end{prob}

\begin{prob}
وښيئ چې $\Nat^\omega$، يعنې د طبيعي عددونو د نامتناهي
تسلسلونو سټ، په راکمونې استدلال !!{nonenumerable} دے۔
\end{prob}

\begin{prob}
فرض کړئ $P$ له مثبتو صحيح عددونو څخه سټ~$\{0\}$ ته د تابعو سټ
دے، او $Q$ له مثبتو صحيح عددونو څخه سټ~$\{0\}$ ته د
\emph{جزوي} تابعو سټ دے۔ وښيئ چې $P$~!!{enumerable} دے او $Q$
نۀ دے۔ (لارښوونه: د $\Bin^\omega$ د شمېرلو مسئله د~$Q$ شمېرلو
ته راکمه کړئ۔)
\end{prob}

\begin{prob}
فرض کړئ $S$ له مثبتو صحيح عددونو څخه سټ~$\{0,1\}$ ته د ټولو
!!{surjective} تابعو سټ دے؛ يعنې $S$ له ټولو !!{surjective}
~$f\colon \PosInt \to \Bin$ څخه جوړ دے۔ وښيئ چې $S$
!!{nonenumerable} دے۔
\textbf{د ژباړې د نسخې سمون (PSSIZ-003):} په انګرېزي نثر کښې
د دوه غړيز سټ نښه له رياضي حالت څخه بهر وه؛ دلته د نورو سټونو
په شان په رياضي حالت کښې ليکل شوې ده۔
\end{prob}

\begin{prob}
وښيئ چې د ټولو حقيقي عددونو سټ~$\Real$ !!{nonenumerable} دے۔
\end{prob}

\end{document}
